\subsection*{Size 28 and size 29}
The lower bound at size 28 is attained. The two exact values are
\[
                            m(28)=m(29)=84.
\]
At cap 84 the only further admissible modulus after 80 is 84 itself. The lower bound for size 29 follows at once: up to cap 83 at most four entries are outside the cap-60 pool, and every 25-element subset of that pool can cover. At size 28 use the cap-80 exclusion above. For positive witnesses write
\[
 P_{84}=P_{60}\cup\{63,70,72,80,84\}.
\]
The set $P_{84}\setminus\{2,3,8,63\}$ attains the minimum at size 28. The set $P_{84}\setminus\{2,4,6\}$ attains it at size 29. Both contain the triple $16,48,80$, so use its binary contraction proved above. On the 7-coordinate at most ten rows occur: $7,14,21,28,35,42,56,63,70,84$. In each seven-bucket partition some bucket is a singleton since $10<14$. For the other buckets take all unions clipped to the singleton mask, intersect across buckets, and retain maximal masks up to translation as before. With all ten rows 41 projected masks on the 360-cycle suffice. The direct residue test returns true on each displayed set and false on every single deletion. Explicit residue assignments in the accompanying files can also be checked by scanning all 5040 integer positions. This establishes the two minima without exhaustive classification of every 28-element optimum.

At size 29 there is in fact just the one extremal set: a 29-set through 84 must contain all five new entries and a noncoverable 24-subset from the cap-60 pool. There are 35 possibilities for that subset. Removing 84 or one of the other new entries already exhibits a smaller cover in most cases; the remaining tests return only the stated extremizer. This count is a finite computational result. At the following size $m(30)\in\{90,96\}$. Indeed a 30-set through 89 already contains a 25-subset from the old pool, and the only addition from 85 through 95 is modulus 90. For the upper bound replace 48 in the 29-element witness by 32 and 96. They are the only rows fixing the new highest binary digit and contract exactly to 48. After a deletion of either, any cover would come from the old rows without 48. Any other single deletion contracts to the corresponding deletion in the 29-set.

\subsection*{Size 30}
The remaining alternative can be eliminated, giving $m(30)=96$. At cap 90 the 30-set needs every one of the six additions $63,70,72,80,84,90$, since otherwise 25 entries come from the old pool. The other 24 entries must form one of the 35 noncoverable old sets. After adjoining the first four additions only six can still be noncoverable; this is the negative list already used at size 29. After adjoining 84 only two remain noncoverable. They are the sets obtained by omitting $2,3,4$ or $2,3,6$ from $P_{84}$. On adjoining 90, neither covers. Test this by the same seven- and ternary-coordinate method: 90 is one additional ordinary 9-divisible row of projected modulus 30 on the 120-cycle. The two whole-set failures exclude every candidate without any further deletion tests. The construction by splitting 48 attains 96 with all 30 moduli distinct; its explicit residues cover period 10080. The finite argument gives the minimum, without counting all 30-element optimizers. At cap 96, keeping both 32 and 96 and contracting them introduces a possibly repeated modulus 48 at the next size.
